% Imported from https://github.com/Luca-Yucheng-Jin/QFT-soln
% Source commit: 0bb3e7066154a8fdb256399a8ecceae720c5e192
\section{David Tong The Standard Model, Sheet 1: The Basics}
\markboth{\thesection\quad TONG THE STANDARD MODEL, SHEET 1}{}

\noindent\textit{These prompts paraphrase David Tong's \emph{The Standard Model},
Example Sheet 1 (January 2026). Problem numbers, starred status, equations,
and guidance follow the
\href{https://davidtong.org/pdfs/teaching/standard-model/rw1.pdf}{official sheet}.}

\subsection{Sheet 1, Problem 1: Lorentz Generators for Weyl Spinors}

Define
\begin{equation}
    \sigma^{\mu\nu}
    =\frac{i}{4}(\sigma^\mu\bar\sigma^\nu
    -\sigma^\nu\bar\sigma^\mu).
\end{equation}
Verify that these matrices obey
\begin{equation}
    [\sigma^{\mu\nu},\sigma^{\rho\sigma}]
    =i\left(\eta^{\nu\rho}\sigma^{\mu\sigma}
    -\eta^{\nu\sigma}\sigma^{\mu\rho}
    +\eta^{\mu\sigma}\sigma^{\nu\rho}
    -\eta^{\mu\rho}\sigma^{\nu\sigma}\right).
\end{equation}
\textit{Hint:} Rewrite $\sigma^{\mu\nu}
=\tfrac i2(\eta^{\mu\nu}-\sigma^\nu\bar\sigma^\mu)$ and
establish
$\sigma^\mu\bar\sigma^\nu
+\sigma^\nu\bar\sigma^\mu=2\eta^{\mu\nu}\mathbf1_2$.
\begin{solution}
    Recall $\{\sigma_i,\sigma_j\} = 2\delta_{ij}$, which follows
    \begin{equation}
        \sigma^\mu\bar\sigma^\nu
+\sigma^\nu\bar\sigma^\mu=2\eta^{\mu\nu}\mathbf1_2.
\label{eq:antisigma}
    \end{equation}
    Hence,
    \begin{equation}
        \sigma^{\mu\nu}
=\tfrac i2(\eta^{\mu\nu}-\sigma^\nu\bar\sigma^\mu),
\label{eq:redef}
    \end{equation}
    which gives
    \begin{equation}
        [\sigma^{\mu\nu},\sigma^{\rho\sigma}] = -\frac{1}{4}(\sigma^\nu\bar\sigma^\mu\sigma^\sigma\bar\sigma^\rho - \sigma^\sigma\bar\sigma^\rho\sigma^\nu\bar\sigma^\mu).
    \end{equation}
    Now the strategy is to permute the indices using \eqref{eq:antisigma} so that the 4 $\sigma$ term eventually cancels. Carrying out the calculation, we find
    \begin{equation}
        [\sigma^{\mu\nu},\sigma^{\rho\sigma}] = -\frac{1}{4}\left(-2\eta^{\rho\nu}\sigma^\sigma \bar \sigma^\mu + 2\eta^{\mu\rho}\sigma^\sigma \bar \sigma^\nu - 2\eta^{\nu\sigma}\sigma^\mu \bar \sigma^\rho + 2\eta^{\mu\sigma}\sigma^\nu \bar \sigma^\rho\right).
    \end{equation}
    Then, plugging \eqref{eq:redef} in, we find
    \begin{equation}
            [\sigma^{\mu\nu},\sigma^{\rho\sigma}]
    =i\left(\eta^{\nu\rho}\sigma^{\mu\sigma}
    -\eta^{\nu\sigma}\sigma^{\mu\rho}
    +\eta^{\mu\sigma}\sigma^{\nu\rho}
    -\eta^{\mu\rho}\sigma^{\nu\sigma}\right).
    \end{equation}
\end{solution}

\subsection{Sheet 1, Problem 2: Raised Weyl Indices and Lorentz Scalars}

Let $S\in SL(2,\mathbb C)$ act on left- and right-handed spinors by
\begin{equation}
    (\psi_L)_\alpha\longmapsto S_\alpha{}^\beta(\psi_L)_\beta,
    \qquad
    (\psi_R)_{\dot\alpha}
    \longmapsto (S^*)_{\dot\alpha}{}^{\dot\beta}
    (\psi_R)_{\dot\beta}.
\end{equation}
The dotted index distinguishes the two representations. Prove the
following statements in the sheet's epsilon convention:
\begin{enumerate}[label=\roman*)]
    \item $(S^{-1})_\beta{}^\alpha
    =\epsilon^{\alpha\gamma}S_\gamma{}^\lambda
    \epsilon_{\lambda\beta}$.
    \item $(\psi_L)^\alpha=\epsilon^{\alpha\beta}(\psi_L)_\beta$
    transforms with $(S^{-1})_\beta{}^\alpha$ on its right.
    \item $\psi_L\chi_L=(\psi_L)^\alpha(\chi_L)_\alpha$
    is Lorentz invariant.
    \item $(\psi_R)^{\dot\alpha}
    =\epsilon^{\dot\alpha\dot\beta}(\psi_R)_{\dot\beta}$
    transforms with $S^{-1\dagger}$.
    \item $\bar\psi_R\psi_L
    =(\psi_R^*)^\alpha(\psi_L)_\alpha$
    is also a Lorentz scalar.
\end{enumerate}
\begin{solution}
    \begin{enumerate}[label=\roman*)]
    \item
    \begin{equation}
        S_\alpha{}^\gamma S_\beta{}^\delta \epsilon_{\gamma \delta} = \begin{pmatrix}
            a &  b\\ c &d
        \end{pmatrix}\begin{pmatrix}
            0 & 1 \\ -1 &0
        \end{pmatrix}
        \begin{pmatrix}
            a & c\\ b &d
        \end{pmatrix} = \begin{pmatrix}
                        0 & ad-bc \\ -ad +bc &0
        \end{pmatrix} = \det S \epsilon_{\alpha \beta}.
    \end{equation}
    Since $S \in SL(2,\mathbb{C})$, $\det S = 1$, we find
    \begin{equation}
        S_\alpha{}^\gamma S_\beta{}^\delta \epsilon_{\gamma \delta} = \epsilon_{\alpha \beta} \implies  \epsilon^{\beta \lambda}S_\alpha{}^\gamma S_\lambda{}^\delta \epsilon_{\gamma \delta} = \delta_\alpha{}^\beta \implies (S^{-1})_\beta{}^\alpha
    =\epsilon^{\alpha\gamma}S_\gamma{}^\lambda
    \epsilon_{\lambda\beta}.
    \end{equation}
    \item \begin{equation}
        (\psi_L)^\alpha = \epsilon^{\alpha \beta } (\psi_L)_\beta \to \epsilon^{\alpha\beta}S_\beta{}^\gamma(\psi_L)_\gamma = \epsilon^{\alpha\beta}S_\beta{}^\gamma\epsilon_{\lambda\gamma}(\psi_L)^\gamma = (S^{-1})_\gamma{}^\alpha (\psi_L)^\gamma.
    \end{equation}
    \item \begin{equation}
        \psi_L \chi_L \to (\psi_L)^\gamma(S^{-1})_\gamma{}^\alpha S_\alpha{}^\gamma(\chi_L)_\beta = \psi_L \chi_L .
    \end{equation}
    \item Completely analogous to the left-handed case, we find
    \begin{equation}
        (\psi_R)^{\dot \alpha} \to (S^{-1\dagger})^{\dot \alpha}{}_{\dot \beta} (\psi_R)^{\dot \beta}.
    \end{equation}
    \item \begin{equation}
        \bar \psi_R \psi_L \to (\psi_R^*)^{\dot \alpha}(S^{-1})_{\dot \alpha}{}^\alpha S_\alpha{}^\beta(\psi_L)_\beta =  \bar \psi_R \psi_L.
    \end{equation}
    \end{enumerate}
\end{solution}

\subsection{Sheet 1, Problem 3: The Right-Handed Lorentz Representation}

In addition to $\sigma^{\mu\nu}$ above, define
\begin{equation}
    \bar\sigma^{\mu\nu}
    =\frac i4(\bar\sigma^\mu\sigma^\nu
    -\bar\sigma^\nu\sigma^\mu).
\end{equation}
Show $\bar\sigma^{\mu\nu}=(\sigma^{\mu\nu})^\dagger$.
If
\begin{equation}
    S=\exp\!\left(-\frac i2\omega_{\mu\nu}
    \sigma^{\mu\nu}\right),
\end{equation}
show that
\begin{equation}
    S^{-1\dagger}
    =\exp\!\left(-\frac i2\omega_{\mu\nu}
    \bar\sigma^{\mu\nu}\right).
\end{equation}
It is enough to verify the infinitesimal transformation and then use
the group law. Relate this to the dotted and undotted components
of a Dirac spinor.
\begin{solution}
    \begin{equation}
        (\sigma^{\mu\nu})^\dagger = -\frac{i}{4}(\bar\sigma^\nu \sigma^\mu - \bar \sigma^\mu \sigma^\nu) = \bar \sigma^{\mu\nu}.
    \end{equation}
    For an infinitesimal transformation, we have
    \begin{equation}
        S^{-1\dagger} = 1 +\left(\frac{i}{2}\omega_{\mu\nu}\sigma^{\mu\nu}\right)^\dagger = 1 - \frac{i}{2}\omega_{\mu\nu}\bar \sigma^{\mu\nu}.
    \end{equation}
    Thereby, the exponential map is given by
    \begin{equation}
            S^{-1\dagger}
    =\exp\!\left(-\frac i2\omega_{\mu\nu}
    \bar\sigma^{\mu\nu}\right).
    \end{equation}
\end{solution}

\subsection{Sheet 1, Problem 4: Dirac and Majorana Mass Mixing}

Two Weyl spinors have a real Dirac mass $M$ and complex Majorana
masses $m_1,m_2$:
\begin{equation}
    \mathcal L_{\rm mass}
    =-M(\bar\psi_L\psi_R+\bar\psi_R\psi_L)
    +\frac{m_1}{2}\psi_L\psi_L
    +\frac{m_1^*}{2}\bar\psi_L\bar\psi_L
    +\frac{m_2}{2}\psi_R\psi_R
    +\frac{m_2^*}{2}\bar\psi_R\bar\psi_R.
\end{equation}
Determine the physical fermion masses.
\begin{solution}
    The mass Lagrangian can be written as
    \begin{equation}
        \mathcal L_{\rm mass} = -\frac{1}{2}\begin{pmatrix}
            \bar \psi _L & \psi_R
        \end{pmatrix}\begin{pmatrix}
            -m_1^* & M\\ M & -m_2^*
        \end{pmatrix}\begin{pmatrix}
            \bar \psi _L \\ \psi_R
        \end{pmatrix} + h.c. \equiv -\frac{1}{2} \chi_i \mathcal{M}_{ij}\chi_j + h.c.,
    \end{equation}
    where $\chi^T \equiv (\bar \psi_L, \psi_R)$.
    The kinetic part of the Lagrangian can be written as
    \begin{equation}
        \bar \psi_L i\bar \sigma^\mu \partial_\mu + \bar\psi_R i\sigma^\mu \partial_\mu \psi_R = \frac{1}{2}(
\bar\chi_i i\bar\sigma^\mu\partial_\mu\chi_i
+
\chi_i i\sigma^\mu\partial_\mu\bar\chi_i).
    \end{equation}
    The equation of motion is given by
    \begin{equation}
        i \sigma^\mu \partial_\mu \chi_i -\mathcal{M}_{ij}^*\bar \chi_j = 0
    \end{equation}
    and its conjugate
    \begin{equation}
        i\bar \sigma^\mu \partial_\mu \bar\chi_i - \mathcal{M}_{ij}\chi_j  = 0.
    \end{equation}
    Applying $i\bar\sigma^\nu  \partial_\nu$ to the first equation, we find
    \begin{equation}
        -\bar\sigma^\nu \sigma^\mu \partial_\mu \partial_\nu \chi_i - \mathcal{M}_{ij}^*i\bar\sigma^\nu \partial_\nu \chi_j =-\bar\sigma^{(\nu} \sigma^{\mu)} \partial_\mu \partial_\nu \chi_i - \mathcal{M}_{ij}^*\mathcal{M}_{jk}\chi_k = 0 \implies (\partial^2 + \mathcal{M}^\dagger\mathcal{M}) \chi = 0.
    \end{equation}
    To progress, let us diagonalise
\begin{equation}
    \mathcal{M}^\dagger \mathcal{M} = U D U^\dagger.
\end{equation}
Then, the equation of motion becomes
\begin{equation}
    (\partial^2 + UDU^\dagger) \chi = 0
    \implies
    (\partial^2 + D)(U^\dagger\chi) = 0.
\end{equation}
Therefore, the physical masses squared are the eigenvalues of
$\mathcal{M}^\dagger\mathcal{M}$. The mass eigenvalues are found to be
\begin{equation}
        m_\pm^2
    =
    \frac{1}{2}\left[
    |m_1|^2+|m_2|^2+2M^2
    \pm
    \sqrt{
    \left(|m_1|^2+|m_2|^2+2M^2\right)^2
    -4|m_1m_2-M^2|^2
    }
    \right].
\end{equation}

\end{solution}

\subsection{Sheet 1, Problem 6*: Discrete Symmetries of Gauge Fields}

Use the chiral-basis matrices
\begin{equation}
    \gamma^\mu=\begin{pmatrix}0&\sigma^\mu\\
    \bar\sigma^\mu&0\end{pmatrix},
    \qquad
    \gamma^5=\begin{pmatrix}\mathbf1&0\\0&-\mathbf1\end{pmatrix},
    \qquad
    \sigma^\mu=(\mathbf1,\sigma^i),\quad
    \bar\sigma^\mu=(\mathbf1,-\sigma^i).
\end{equation}
For a Dirac fermion coupled to $U(1)$, the gauge-field equation is
$\partial_\nu F^{\mu\nu}=e\bar\psi\gamma^\mu\psi$.
Starting from the spinor actions of $C,P,T$, work out how $A_\mu$
must transform so that this equation remains invariant. Then determine
how the Maxwell theta term
\begin{equation}
    S_\theta=\int d^4x\,F_{\mu\nu}{}^\star F^{\mu\nu}
\end{equation}
transforms under each of $C,P,T$.
\begin{solution}
    Under parity, we have
    \begin{equation}
        P: \quad \psi(t,\mathbf{x}) \to \gamma^0\psi(t,-\mathbf{x}), \quad \bar \psi \gamma^\mu \psi(t,\mathbf{x}) \to   \bar \psi\gamma^\mu \psi(t,-\mathbf{x}),
    \end{equation}
    which implies
    \begin{equation}
        P: \quad \partial_\mu F^{\mu 0}(t,\mathbf{x}) \to \partial_\mu F^{\mu 0}(t,-\mathbf{x}), \quad \partial_\mu F^{\mu i}(t,\mathbf{x}) \to - \partial_\mu F^{\mu i}(t,-\mathbf{x}).
    \end{equation}
    This follows
    \begin{equation}
        P: \quad \partial_i\partial^i A^0 - \partial_i \partial^0 A^i \to  \partial_i\partial^i A'^0 + \partial_i \partial^0 A'^i \implies P:\quad A^0(t,\mathbf{x}) \to A^0(t,-\mathbf{x}), \quad A^i(t,\mathbf{x}) \to -A^i(t,-\mathbf{x}).
    \end{equation}
    One can check that the spatial components also transform as desired.

    Under charge conjugation, we have
    \begin{equation}
        C: \quad \psi(t,\mathbf{x}) \to \pm i \gamma^2 \psi^* (t,\mathbf{x}), \quad \bar \psi \gamma^\mu \psi(t,\mathbf{x}) \to -\psi^T\gamma^2 \gamma^0 \gamma^\mu \gamma^2 \psi^* (t,\mathbf{x}) = -\bar \psi \gamma^\mu \psi (t,\mathbf{x}),
    \end{equation}
    where we have used $\gamma^2 \gamma^\mu \gamma^2 = - \gamma^2$.
    Therefore, we must have
    \begin{equation}
        C: \quad A^\mu (t,\mathbf{x}) \to -A^\mu (t,\mathbf{x}) .
    \end{equation}

    Under time reversal, we have
    \begin{equation}
    \begin{aligned}
        T: \quad \psi(t,\mathbf{x})  \to \gamma^1\gamma^3 \psi^*(-t,\mathbf{x}), \quad \bar \psi\gamma^0 \psi (t,\mathbf{x}) &\to   \psi^T (\gamma^0)^*\psi^*(-t,\mathbf{x}) =  \bar \psi \gamma^0\psi(-t,\mathbf{x}),\\  \bar \psi\gamma^i \psi (t,\mathbf{x}) &\to   -\psi^T (\gamma^i)^*\psi^*(-t,\mathbf{x}) =  -\bar \psi \gamma^i\psi(-t,\mathbf{x}).
    \end{aligned}
    \end{equation}
    The LHS of the gauge-field equation transforms as
    \begin{equation}
        T: \quad \partial_\mu F^{\mu0}(t,\mathbf{x}) \to (\partial_j \partial^j A'^0  + \partial_j \partial^0 A'^j)(-t,\mathbf{x}), \quad \partial_\mu F^{\mu i}(t,\mathbf{x}) \to (\partial^2A'^i + \partial_0 \partial^j A'^0 - \partial_i \partial^jA'^j)(-t,\mathbf{x}).
    \end{equation}
    For the theory to be invariant under time reversal, we must have
    \begin{equation}
        T: \quad A^0(t,\mathbf{x}) \to A^0(-t,\mathbf{x}), \quad A^i(t,\mathbf{x}) \to - A^i(-t,\mathbf{x}).
    \end{equation}

    The theta term is proportional to
    \begin{equation}
        \int d^4x \mathbf{E}\cdot \mathbf{B},
    \end{equation}
    whereas the fields transform as
    \begin{equation}
        P:\quad \mathbf{E}(t,\mathbf{x}) \to -\mathbf{E}(t,-\mathbf{x}), \quad \mathbf{B}(t,\mathbf{x}) \to \mathbf{B}(t,-\mathbf{x})
    \end{equation}
    \begin{equation}
        C:\quad \mathbf{E}(t,\mathbf{x}) \to -\mathbf{E}(t,\mathbf{x}), \quad \mathbf{B}(t,\mathbf{x}) \to - \mathbf{B}(t,\mathbf{x})
    \end{equation}
    \begin{equation}
        P:\quad \mathbf{E}(t,\mathbf{x}) \to \mathbf{E}(-t,\mathbf{x}), \quad \mathbf{B}(t,\mathbf{x}) \to - \mathbf{B}(-t,\mathbf{x}).
    \end{equation}
    Therefore, the theta term is invariant under C, and odd under P and T.
\end{solution}

\subsection{Sheet 1, Problem 7: Parity with Scalar and Pseudoscalar Masses}

Prove that $\bar\psi\psi$ and $i\bar\psi\gamma^5\psi$ are real.
For
\begin{equation}
    \mathcal L_{\rm mass}
    =m_1\bar\psi\psi+im_2\bar\psi\gamma^5\psi,
\end{equation}
rewrite the mass term in left- and right-handed Weyl fields using the
chiral gamma matrices. Find an action of parity on $\psi_L,\psi_R$
that preserves the theory, and express it on the Dirac spinor.
\begin{solution}
    \begin{equation}
        (\bar\psi \psi)^* = \psi^T\gamma^0 \psi^* = \bar \psi \psi,
    \end{equation}
    \begin{equation}
        (i\bar \psi \gamma^5 \psi)^* = -i\psi^T \gamma^0 \gamma^5 \psi^* =  - i\psi^\dagger \gamma^5\gamma^0 \psi = i \bar \psi \gamma^5 \psi.
    \end{equation}
    In Weyl spinors, the mass terms are
    \begin{equation}
        m_1(\psi_L^\dagger\psi_R + \psi_R^\dagger \psi_L) + im_2(\psi_L^\dagger\psi_R - \psi_R^\dagger \psi_L).
    \end{equation}
    For the whole Lagrangian to be invariant under parity, we must have
    \begin{equation}
        P :\quad \psi_L(t,\mathbf{x}) \to e^{i\theta}\psi _R (t,-\mathbf{x}), \quad\psi_R(t,\mathbf{x}) \to e^{i\theta}\psi _L (t,-\mathbf{x}),
    \end{equation}
    and hence
    \begin{equation}
        P :\quad \psi(t,\mathbf{x}) \to e^{i\theta}\gamma^0\psi (t,-\mathbf{x}).
    \end{equation}
\end{solution}

\subsection{Sheet 1, Problem 8: Parity and Mass in $2+1$ Dimensions}

In signature $(+,-,-)$ use imaginary matrices
$\gamma^\mu=(\sigma^2,i\sigma^1,i\sigma^3)$ and the action
\begin{equation}
    S=-\int d^dx\left(i\bar\psi\gamma^\mu\partial_\mu\psi
    -M\bar\psi\psi\right).
\end{equation}
Explain why parity reverses only $x^1$ and leaves $x^0,x^2$
unchanged, instead of sending the whole spatial vector to its
negative. For $M=0$, find parity, charge-conjugation, and
time-reversal actions that preserve $S$. Which are broken by $M\ne0$?
\begin{solution}
    We wish to use the parity action to enlarge $SO(1,2)$ to $O(1,2)$. To do that, we need the determinant of parity to be negative. Therefore, this can only be achieved with one coordinate flipped in two spatial dimensions.

    For parity, the massless Lagrangian transform as
    \begin{equation}
        P:\quad i\bar \psi \sigma^\mu \partial_\mu \psi \to i\psi^\dagger P^\dagger\gamma^0(\gamma^0\partial_0 -  \gamma^1 \partial_1 + \gamma^2 \partial_2)P\psi.
    \end{equation}
    Therefore, by choosing $P: \psi \to \gamma^0 \psi$, the massless theory remains invariant.

    For charge conjugation, the Dirac field transforms as
    \begin{equation}
        C: \quad \psi \to C\psi^*,
    \end{equation}
    while the kinetic term transforms as
    \begin{equation}
        C: \quad i\bar\psi\gamma^\mu \partial_\mu \psi \to i \psi^TC^\dagger \gamma^0 \gamma^\mu C\partial_\mu\psi^* = -i(\partial_\mu \psi^\dagger)C^T (\gamma^\mu)^T (\gamma^0)^TC^*\psi = i\psi^\dagger C^T\gamma^\mu \gamma^0 C^* \psi,
    \end{equation}
    where we have used $(\gamma^\mu)^T = -\gamma^\mu$.
    Therefore, we must have $C:  \psi \to e^{i\alpha}\psi^*$.

    For time reversal, we have
    \begin{equation}
        T: \quad i \to -i, \quad \psi \to T \psi^*,
    \end{equation}
    and hence
    \begin{equation}
    \begin{aligned}
        T: \quad i\bar \psi \gamma^\mu \partial_\mu \psi &\to -i \psi^T T^\dagger \gamma^0 (-\gamma^0 T \partial_0 -\gamma^1 T \partial_1-\gamma^2 T \partial_2)\psi^*\\ &= -i\psi^\dagger T^T (-\gamma^0 \gamma^0 T^* \partial_0 - \gamma^1 \gamma^0 T^* \partial_1 - \gamma^2 \gamma^0 T^* \partial_2) \psi.
    \end{aligned}
    \end{equation}
    For the kinetic term to be invariant, we must have $T:\psi \to \gamma^1\gamma^2\psi^* = \gamma^0 \psi^*$.

    Now it is clear that the mass term is even under C, and odd under T and P .
\end{solution}

